\documentclass[11pt,reqno]{amsart}

\usepackage[T1]{fontenc}
\usepackage{lmodern}
\usepackage{microtype}
\usepackage{mathtools}
\usepackage{amssymb}
\usepackage{amsthm}
\usepackage{booktabs}
\usepackage{array}
\usepackage{xcolor}
\usepackage[margin=1.15in]{geometry}
\usepackage[colorlinks=true,linkcolor=blue!55!black,citecolor=blue!55!black,
  urlcolor=blue!55!black]{hyperref}
\usepackage[capitalize,noabbrev]{cleveref}

\allowdisplaybreaks
\raggedbottom

\newtheorem{theorem}{Theorem}[section]
\newtheorem{proposition}[theorem]{Proposition}
\newtheorem{lemma}[theorem]{Lemma}
\newtheorem{corollary}[theorem]{Corollary}
\theoremstyle{definition}
\newtheorem{definition}[theorem]{Definition}
\newtheorem{remark}[theorem]{Remark}

\crefname{equation}{equation}{equations}
\crefname{definition}{definition}{definitions}
\crefname{theorem}{theorem}{theorems}
\crefname{lemma}{lemma}{lemmas}
\crefname{proposition}{proposition}{propositions}
\crefname{remark}{remark}{remarks}
\crefname{section}{section}{sections}

\DeclareMathOperator{\rank}{rank}
\DeclareMathOperator{\supp}{supp}
\newcommand{\R}{\mathbb{R}}
\newcommand{\M}{\mathbb{M}}
\newcommand{\V}{\mathcal{V}}
\newcommand{\Kprod}{\mathcal{K}}
\newcommand{\trans}{\mathsf{T}}
\newcommand{\dd}{\,\mathrm{d}}
\newcommand{\abs}[1]{\lvert#1\rvert}
\newcommand{\norm}[1]{\lVert#1\rVert}
\newcommand{\set}[1]{\left\{#1\right\}}
\newcommand{\restr}[1]{\big|_{#1}}

\title[An explicit constitutive $T_4$-configuration]
{An Explicit Nondegenerate $T_4$-Configuration\\
in the Constitutive Set of a Genuinely Nonlinear\\
$2\times2$ System of Conservation Laws}

\date{\today}
\subjclass[2020]{35L65, 35L67, 49Q20}
\keywords{systems of conservation laws, strict hyperbolicity, genuine
nonlinearity, convex entropy, $T_N$-configuration, constitutive set}

\begin{document}

\begin{abstract}
We give a constructive example of a smooth flux for a one-dimensional
$2\times2$ system of conservation laws which is strictly hyperbolic and
genuinely nonlinear on a convex open state domain and which admits a strictly
convex entropy.  The associated entropy-augmented constitutive set in
$\R^{3\times2}$ contains an ordered $T_4$-configuration.  The configuration is
nondegenerate in the stronger sense that no two of its four wells are
rank-one connected.  All matrices and all finite interpolation data are
rational and explicit.  The flux is obtained from a scalar potential.  A
one-dimensional Hermite interpolant places the four prescribed wells, while
a transverse smooth connector enforces genuine nonlinearity of both
characteristic families.  We include the exact rank calculations, the full
connector construction, and the passage from the central state segment to a
convex open strip.
\end{abstract}

\maketitle

\section{Introduction and main result}

Let $\V\subset\R^2$ be a nonempty convex open set and consider the system of
conservation laws
\begin{equation}
  U_t+F(U)_x=0,
  \qquad U\colon\R\times[0,\infty)\longrightarrow\V,
  \label{eq:system}
\end{equation}
with a smooth flux $F=(F_1,F_2)\colon\V\to\R^2$.  We regard differentials of
scalar functions as row covectors.  Thus a pair of smooth scalar functions
$(\eta,q)$ is an entropy--entropy-flux pair for \cref{eq:system} when
\begin{equation}
  Dq(U)=D\eta(U)\,DF(U),
  \qquad U\in\V.
  \label{eq:entropy-compatibility}
\end{equation}
The entropy is strictly convex when $D^2\eta(U)$ is positive definite for
every $U\in\V$.

The system is strictly hyperbolic when $DF(U)$ has two distinct real
eigenvalues
\begin{equation}
  \lambda_1(U)<\lambda_2(U),
  \qquad U\in\V.
\end{equation}
If $r_k(U)$ is a corresponding right eigenvector, the $k$th characteristic
field is genuinely nonlinear when
\begin{equation}
  D\lambda_k(U)[r_k(U)]\neq0,
  \qquad U\in\V,\qquad k=1,2.
  \label{eq:genuine-nonlinearity}
\end{equation}

For $U=(u_1,u_2)\in\V$, define
\begin{equation}
  G_{F,\eta,q}(U)
  :=
  \begin{pmatrix}
    u_1 & -F_1(U)\\
    u_2 & -F_2(U)\\
    \eta(U) & -q(U)
  \end{pmatrix}.
  \label{eq:G-map}
\end{equation}
The entropy-augmented constitutive set is
\begin{equation}
  K_{F,\eta,q}
  :=G_{F,\eta,q}(\V)
  \subset\R^{3\times2}.
  \label{eq:constitutive-set}
\end{equation}
For a fixed positive integer $N$, it is useful to distinguish the single-well
set in \cref{eq:constitutive-set} from its $N$-fold product:
\begin{equation}
  \Kprod_{F,\eta,q}^{(N)}
  :=\bigl(K_{F,\eta,q}\bigr)^N
  \subset\bigl(\R^{3\times2}\bigr)^N.
  \label{eq:product-constitutive-set}
\end{equation}

Our main result is the following explicit affirmative construction.

\begin{theorem}[A nondegenerate constitutive $T_4$]
  \label{thm:main}
There exist a number $\delta>0$, a convex open rectangle
\begin{equation}
  \V=(-8,8)\times(-\delta,\delta),
  \label{eq:state-domain}
\end{equation}
smooth functions
\begin{equation}
  F\colon\V\to\R^2,
  \qquad
  \eta,q\colon\V\to\R,
\end{equation}
and four states $U_1,\ldots,U_4\in\V$ with the following properties:
\begin{enumerate}
  \item $D^2\eta=I_2$, so $\eta$ is uniformly strictly convex;
  \item $(\eta,q)$ satisfies \cref{eq:entropy-compatibility};
  \item the system \cref{eq:system} is strictly hyperbolic throughout $\V$;
  \item both characteristic fields are genuinely nonlinear throughout $\V$;
  \item the ordered tuple
  \begin{equation}
    \bigl(G_{F,\eta,q}(U_1),\ldots,G_{F,\eta,q}(U_4)\bigr)
    \in\Kprod_{F,\eta,q}^{(4)}
  \end{equation}
  is an ordered $T_4$-configuration;
  \item no two distinct members of this configuration are rank-one
  connected.
\end{enumerate}
In fact, one may take
\begin{equation}
  U_1=(-7,0),\qquad U_2=(-1,0),\qquad
  U_3=(1,0),\qquad U_4=(7,0),
  \label{eq:four-states}
\end{equation}
and the flux may be chosen to be a gradient, $F=\nabla\Phi$, for a smooth
scalar potential $\Phi$.
\end{theorem}

The proof separates into two parts.  First, we give an exact rational
$T_4$-configuration and determine the finite jets that a potential must
realize at the states in \cref{eq:four-states}.  Second, we construct a smooth
potential realizing those jets while preserving strict hyperbolicity and
genuine nonlinearity on a thin convex strip.

\section{Ordered \texorpdfstring{$T_N$}{TN}-configurations and potential fluxes}

Set
\begin{equation}
  \M:=\R^{3\times2},
  \qquad
  \mathcal R_1:=\set{C\in\M:\rank C=1}.
\end{equation}
Thus $\mathcal R_1$ is the manifold of nonzero rank-one matrices in $\M$.

\begin{definition}[Ordered $T_N$-configuration]
  \label{def:TN}
Fix $N\ge4$.  An ordered $N$-tuple
\begin{equation}
  \mathbf X=(X_1,\ldots,X_N)\in\M^N
\end{equation}
is an ordered $T_N$-configuration if there exist
\begin{equation}
  P\in\M,
  \qquad C_i\in\mathcal R_1,
  \qquad \kappa_i>1,
  \qquad i=1,\ldots,N,
  \label{eq:TN-parameters}
\end{equation}
such that
\begin{equation}
  \sum_{i=1}^N C_i=0
  \label{eq:TN-closure}
\end{equation}
and
\begin{equation}
  X_i=P+\sum_{j=1}^{i-1}C_j+\kappa_iC_i,
  \qquad i=1,\ldots,N.
  \label{eq:TN-parametrization}
\end{equation}
The tuple
\begin{equation}
  \bigl(P,C_1,\ldots,C_N,\kappa_1,\ldots,\kappa_N\bigr)
\end{equation}
is called an ordered $T_N$-parameterization of $\mathbf X$.
\end{definition}

\begin{definition}[Nondegeneracy]
  \label{def:nondegenerate}
An ordered $T_N$-configuration in $\R^{3\times2}$ is called nondegenerate if
\begin{equation}
  \rank(X_i-X_j)=2,
  \qquad i\neq j.
  \label{eq:no-rank-one-connections}
\end{equation}
Because a $3\times2$ matrix has rank at most two, condition
\cref{eq:no-rank-one-connections} is precisely the assertion that no two
distinct wells are rank-one connected.
\end{definition}

The gradient-flux ansatz makes entropy compatibility automatic.

\begin{lemma}[Potential-generated entropy pair]
  \label{lem:potential-entropy}
Let $\Phi$ be smooth on an open subset of $\R^2$ and define
\begin{equation}
  F(U):=\nabla\Phi(U),
  \qquad
  \eta(U):=\frac12\abs{U}^2,
  \qquad
  q(U):=U\cdot\nabla\Phi(U)-\Phi(U).
  \label{eq:potential-ansatz}
\end{equation}
Then $D^2\eta=I_2$ and $(\eta,q)$ satisfies
\cref{eq:entropy-compatibility}.
\end{lemma}

\begin{proof}
For $v\in\R^2$, the product rule gives
\begin{align}
  Dq(U)[v]
  &=v\cdot\nabla\Phi(U)
    +U\cdot D^2\Phi(U)v-D\Phi(U)[v]\\
  &=U\cdot D^2\Phi(U)v,
\end{align}
because $D\Phi(U)[v]=\nabla\Phi(U)\cdot v$.  Since
\begin{equation}
  D\eta(U)=U^{\trans},
  \qquad
  DF(U)=D^2\Phi(U),
\end{equation}
we obtain $Dq(U)=D\eta(U)DF(U)$.  The identity
$D^2\eta=I_2$ is immediate.
\end{proof}

We shall also use the following standard eigenvalue differentiation formula,
for which we include the proof to fix all conventions.

\begin{lemma}[Directional derivative of a simple Hessian eigenvalue]
  \label{lem:eigenvalue-derivative}
Let $\Phi\in C^3(\Omega)$, let $D^2\Phi(U)$ have a simple eigenvalue
$\lambda_k(U)$, and let $r_k(U)$ be a differentiable unit eigenvector.  Then,
for every $v\in\R^2$,
\begin{equation}
  D\lambda_k(U)[v]
  =D^3\Phi(U)[v,r_k(U),r_k(U)].
  \label{eq:eigenvalue-derivative}
\end{equation}
In particular,
\begin{equation}
  D\lambda_k(U)[r_k(U)]
  =D^3\Phi(U)[r_k(U),r_k(U),r_k(U)].
  \label{eq:genuine-cubic}
\end{equation}
\end{lemma}

\begin{proof}
Differentiate
\begin{equation}
  D^2\Phi(U)r_k(U)=\lambda_k(U)r_k(U)
\end{equation}
in the direction $v$ and take the Euclidean inner product with $r_k(U)$.
The terms containing $Dr_k(U)[v]$ cancel because $D^2\Phi(U)$ is symmetric
and $D^2\Phi(U)r_k(U)=\lambda_k(U)r_k(U)$.  Since $\abs{r_k(U)}=1$, the
remaining identity is exactly \cref{eq:eigenvalue-derivative}.  Taking
$v=r_k(U)$ gives \cref{eq:genuine-cubic}.
\end{proof}

\section{An exact rational \texorpdfstring{$T_4$}{T4} certificate}
\label{sec:algebraic-T4}

We now give the algebraic core of the construction.  Take
\begin{equation}
  \kappa_1=\kappa_2=\kappa_3=\kappa_4=2
  \label{eq:kappas}
\end{equation}
and
\begin{equation}
  P=
  \begin{pmatrix}
    \frac{17}{5}&0\\
    0&0\\
    \frac{149}{10}&0
  \end{pmatrix}.
  \label{eq:P}
\end{equation}
Define the four increments by
\begin{align}
  C_1&=
  \begin{pmatrix}
    -\frac{26}{5}&-\frac{26}{5}\\
    0&0\\
    \frac{24}{5}&\frac{24}{5}
  \end{pmatrix},
  &
  C_2&=
  \begin{pmatrix}
    \frac25&-\frac45\\
    0&0\\
    -\frac{48}{5}&\frac{96}{5}
  \end{pmatrix},
  \label{eq:C12}\\[1ex]
  C_3&=
  \begin{pmatrix}
    \frac65&6\\
    0&0\\
    -\frac{24}{5}&-24
  \end{pmatrix},
  &
  C_4&=
  \begin{pmatrix}
    \frac{18}{5}&0\\
    0&0\\
    \frac{48}{5}&0
  \end{pmatrix}.
  \label{eq:C34}
\end{align}
Their rank-one structure is transparent from the outer-product
factorizations
\begin{align}
  C_1&=
  \begin{pmatrix}-\frac{26}{5}\\[1mm]0\\[1mm]\frac{24}{5}\end{pmatrix}
  \begin{pmatrix}1&1\end{pmatrix},
  &
  C_2&=
  \begin{pmatrix}\frac25\\[1mm]0\\[1mm]-\frac{48}{5}\end{pmatrix}
  \begin{pmatrix}1&-2\end{pmatrix},
  \label{eq:factor-C12}\\[1ex]
  C_3&=
  \begin{pmatrix}\frac65\\[1mm]0\\[1mm]-\frac{24}{5}\end{pmatrix}
  \begin{pmatrix}1&5\end{pmatrix},
  &
  C_4&=
  \begin{pmatrix}\frac{18}{5}\\[1mm]0\\[1mm]\frac{48}{5}\end{pmatrix}
  \begin{pmatrix}1&0\end{pmatrix}.
  \label{eq:factor-C34}
\end{align}
Every column factor and every row factor in
\cref{eq:factor-C12,eq:factor-C34} is nonzero.  Hence
$C_i\in\mathcal R_1$ for every $i$.  Direct summation gives
\begin{equation}
  \sum_{i=1}^4C_i
  =
  \begin{pmatrix}
    \frac{-26+2+6+18}{5}&\frac{-26-4+30}{5}\\
    0&0\\
    \frac{24-48-24+48}{5}&\frac{24+96-120}{5}
  \end{pmatrix}
  =0_{3\times2},
  \label{eq:closure-check}
\end{equation}
which verifies \cref{eq:TN-closure}.

Using \cref{eq:TN-parametrization} with the order displayed above and
$\kappa_i=2$, we obtain
\begin{align}
  X_1&=P+2C_1
  =
  \begin{pmatrix}
    -7&-\frac{52}{5}\\
    0&0\\
    \frac{49}{2}&\frac{48}{5}
  \end{pmatrix},
  &
  X_2&=P+C_1+2C_2
  =
  \begin{pmatrix}
    -1&-\frac{34}{5}\\
    0&0\\
    \frac12&\frac{216}{5}
  \end{pmatrix},
  \label{eq:X12}\\[1ex]
  X_3&=P+C_1+C_2+2C_3
  =
  \begin{pmatrix}
    1&6\\
    0&0\\
    \frac12&-24
  \end{pmatrix},
  &
  X_4&=P+C_1+C_2+C_3+2C_4
  =
  \begin{pmatrix}
    7&0\\
    0&0\\
    \frac{49}{2}&0
  \end{pmatrix}.
  \label{eq:X34}
\end{align}
Thus $(X_1,X_2,X_3,X_4)$ is an ordered $T_4$-configuration.  It remains at
this algebraic stage to verify the stronger pairwise nondegeneracy property.

For $1\le i<j\le4$, let
\begin{equation}
  \Delta_{ij}
  :=\det\left((X_i-X_j)_{\{1,3\}\times\{1,2\}}\right),
  \label{eq:minor-definition}
\end{equation}
where we retain the first and third rows and both columns.  The first three
submatrices and determinants are
\begin{equation}
\renewcommand{\arraystretch}{1.2}
\begin{array}{c|c|c}
  (i,j)&(X_i-X_j)_{\{1,3\}\times\{1,2\}}&\Delta_{ij}\\
  \hline
  (1,2)&\begin{pmatrix}-6&-\frac{18}{5}\\24&-\frac{168}{5}\end{pmatrix}
       &288\\
  (1,3)&\begin{pmatrix}-8&-\frac{82}{5}\\24&\frac{168}{5}\end{pmatrix}
       &\frac{624}{5}\\
  (1,4)&\begin{pmatrix}-14&-\frac{52}{5}\\0&\frac{48}{5}\end{pmatrix}
       &-\frac{672}{5}
\end{array}
\label{eq:minor-table-first}
\end{equation}
and the remaining three are
\begin{equation}
\renewcommand{\arraystretch}{1.2}
\begin{array}{c|c|c}
  (i,j)&(X_i-X_j)_{\{1,3\}\times\{1,2\}}&\Delta_{ij}\\
  \hline
  (2,3)&\begin{pmatrix}-2&-\frac{64}{5}\\0&\frac{336}{5}\end{pmatrix}
       &-\frac{672}{5}\\
  (2,4)&\begin{pmatrix}-8&-\frac{34}{5}\\-24&\frac{216}{5}\end{pmatrix}
       &-\frac{2544}{5}\\
  (3,4)&\begin{pmatrix}-6&6\\-24&-24\end{pmatrix}
       &288
\end{array}
\label{eq:minor-table-second}
\end{equation}
Every determinant in
\cref{eq:minor-table-first,eq:minor-table-second} is nonzero.  Therefore
$\rank(X_i-X_j)\ge2$ for $i\neq j$.  Since these are $3\times2$ matrices,
their ranks are at most two, and hence
\begin{equation}
  \rank(X_i-X_j)=2,
  \qquad i\neq j.
  \label{eq:rank-two-conclusion}
\end{equation}

\section{Reduction to finite potential jets}
\label{sec:jets}

We now realize the matrices in \cref{eq:X12,eq:X34} as constitutive wells.
Use the potential ansatz from \cref{lem:potential-entropy}.  Write a state as
$U=(x,y)$ and let
\begin{equation}
  \eta(x,y)=\frac{x^2+y^2}{2}.
\end{equation}
Suppose that $\Phi_y(x,0)=0$ and set
\begin{equation}
  \psi(x):=\Phi(x,0).
\end{equation}
At a state $(x,0)$ we then have
\begin{equation}
  F(x,0)=\bigl(\psi'(x),0\bigr),
  \qquad
  q(x,0)=x\psi'(x)-\psi(x),
  \label{eq:axis-flux-q}
\end{equation}
and consequently
\begin{equation}
  G_{F,\eta,q}(x,0)
  =
  \begin{pmatrix}
    x&-\psi'(x)\\
    0&0\\
    \frac{x^2}{2}&\psi(x)-x\psi'(x)
  \end{pmatrix}.
  \label{eq:axis-G}
\end{equation}

Let
\begin{equation}
  \xi_1=-7,
  \qquad \xi_2=-1,
  \qquad \xi_3=1,
  \qquad \xi_4=7.
\end{equation}
Comparing \cref{eq:axis-G} with \cref{eq:X12,eq:X34}, one sees that it is
enough to impose $\Phi_y(\xi_i,0)=0$ together with the following scalar
Hermite data:
\begin{equation}
\renewcommand{\arraystretch}{1.25}
\begin{array}{c|rrrr}
  x&-7&-1&1&7\\
  \hline
  \psi(x)&-\frac{316}{5}&\frac{182}{5}&-30&0\\[1mm]
  \psi'(x)&\frac{52}{5}&\frac{34}{5}&-6&0
\end{array}.
\label{eq:Hermite-data}
\end{equation}
For example, at $x=-7$ the upper-right entry of $X_1$ forces
$\psi'(-7)=52/5$, and then its lower-right entry forces
\begin{equation}
  \psi(-7)-(-7)\psi'(-7)=\frac{48}{5},
\end{equation}
which gives $\psi(-7)=-316/5$.  The other three columns of
\cref{eq:Hermite-data} follow by the same calculation.

\section{The one-dimensional interpolant and the transverse connector}
\label{sec:connector}

\subsection{The Hermite polynomial}

Define the degree-seven polynomial
\begin{align}
  \psi_H(x)
  :={}&
  \frac{33263}{5120}
  -\frac{4085033}{80640}x
  -\frac{121651}{35840}x^2
  +\frac{23868737}{1317120}x^3\notag\\
  &+\frac{3531}{35840}x^4
  -\frac{176725}{263424}x^5
  -\frac{33}{35840}x^6
  +\frac{26129}{3951360}x^7.
  \label{eq:psi-H}
\end{align}
This is the unique polynomial of degree at most seven satisfying the eight
conditions in \cref{eq:Hermite-data}.  Introduce
\begin{equation}
  R(x):=(x^2-1)(x^2-49)
  \label{eq:R-polynomial}
\end{equation}
and set
\begin{equation}
  \psi(x):=\psi_H(x)-xR(x)^2.
  \label{eq:psi-definition}
\end{equation}
At each $\xi_i\in\{-7,-1,1,7\}$ one has $R(\xi_i)=0$.  Therefore
\begin{equation}
  \bigl(xR(x)^2\bigr)\restr{x=\xi_i}=0
\end{equation}
and, since
\begin{equation}
  \frac{\dd}{\dd x}\bigl(xR(x)^2\bigr)=R(x)^2+2xR(x)R'(x),
\end{equation}
its first derivative also vanishes at every $\xi_i$.  Hence $\psi$ and
$\psi_H$ have the same values and first derivatives at the four nodes.  A
direct exact evaluation gives
\begin{equation}
\renewcommand{\arraystretch}{1.25}
\begin{array}{c|rrrr}
  x&-7&-1&1&7\\
  \hline
  \psi_H(x)=\psi(x)&-\frac{316}{5}&\frac{182}{5}&-30&0\\[1mm]
  \psi_H'(x)=\psi'(x)&\frac{52}{5}&\frac{34}{5}&-6&0,
\end{array}
\label{eq:Hermite-verification}
\end{equation}
which verifies \cref{eq:Hermite-data}.

The correction in \cref{eq:psi-definition} is chosen to create regions in
which the third derivative is positive without altering the required jets.
Indeed, exact differentiation yields
\begin{equation}
  \begin{aligned}
    \psi'''(-5)&=\frac{455187556553}{329280}>0,
    &\qquad
    \psi'''(0)&=\frac{6477756737}{219520}>0,\\
    \psi'''(5)&=\frac{455186246783}{329280}>0.
  \end{aligned}
  \label{eq:positive-third-derivatives}
\end{equation}
By continuity, there are pairwise disjoint open intervals
\begin{equation}
  J_{-5}\Subset(-7,-1),
  \qquad
  J_0\Subset(-1,1),
  \qquad
  J_5\Subset(1,7),
  \label{eq:J-intervals}
\end{equation}
centered at $-5,0,5$, respectively, such that
\begin{equation}
  \overline{J_{-5}}\cup\overline{J_0}\cup\overline{J_5}
  \subset\set{x\in\R:\psi'''(x)>0}.
  \label{eq:J-positive}
\end{equation}

\subsection{A smooth connector with four prescribed zeros}

For completeness, the bump functions used below may be chosen from the
standard profile
\begin{equation}
  \vartheta(t):=
  \begin{cases}
    \exp\!\bigl(-1/(1-t^2)\bigr),&\abs{t}<1,\\
    0,&\abs{t}\ge1.
  \end{cases}
  \label{eq:standard-bump}
\end{equation}
Choose $\rho_a>0$ so small that
$[a-\rho_a,a+\rho_a]\subset J_a$ and take
\begin{equation}
  \alpha_a(x):=\vartheta\!\left(\frac{x-a}{\rho_a}\right),
  \qquad a\in\{-5,0,5\}.
\end{equation}
Thus $\alpha_a\in C_c^\infty(J_a)$ is nonnegative and nonzero.  Normalize
these bumps by
\begin{equation}
  \beta_{-5}:=
  \frac{6\alpha_{-5}}{\int_{\R}\alpha_{-5}(t)\dd t},
  \qquad
  \beta_0:=
  \frac{2\alpha_0}{\int_{\R}\alpha_0(t)\dd t},
  \qquad
  \beta_5:=
  \frac{6\alpha_5}{\int_{\R}\alpha_5(t)\dd t}.
  \label{eq:normalized-bumps}
\end{equation}
Define
\begin{equation}
  g(x):=-1+\beta_{-5}(x)+\beta_0(x)+\beta_5(x),
  \qquad
  h(x):=\int_{-7}^xg(t)\dd t.
  \label{eq:g-h}
\end{equation}
Since the three supports lie in disjoint intervals, we have
\begin{align}
  \int_{-7}^{-1}g(t)\dd t&=-6+\int_{\R}\beta_{-5}(t)\dd t=0,
  \label{eq:g-integral-left}\\
  \int_{-1}^{1}g(t)\dd t&=-2+\int_{\R}\beta_0(t)\dd t=0,
  \label{eq:g-integral-middle}\\
  \int_{1}^{7}g(t)\dd t&=-6+\int_{\R}\beta_5(t)\dd t=0.
  \label{eq:g-integral-right}
\end{align}
It follows successively from
\cref{eq:g-h,eq:g-integral-left,eq:g-integral-middle,eq:g-integral-right} that
\begin{equation}
  h(-7)=h(-1)=h(1)=h(7)=0.
  \label{eq:h-zeros}
\end{equation}

Let
\begin{equation}
  I:=[-8,8],
  \qquad
  E:=\set{x\in I:g(x)\ge0}.
  \label{eq:I-E}
\end{equation}
Outside the union of the three bump supports, $g=-1$.  Consequently,
\begin{equation}
  E\subset
  \bigcup_{a\in\{-5,0,5\}}\supp\beta_a
  \Subset\set{x\in\R:\psi'''(x)>0}.
  \label{eq:E-positive-region}
\end{equation}
The set $E$ is compact.  It is also nonempty: on each of the intervals of
length $6,2,6$, respectively, the corresponding $\beta_a$ has integral equal
to that length and vanishes near both endpoints, so it must exceed $1$
somewhere.  In particular,
\begin{equation}
  \mu:=\min_{x\in E}\psi'''(x)>0.
  \label{eq:mu}
\end{equation}

\subsection{Choice of the potential on the central segment}

Choose the concrete constant
\begin{equation}
  L:=2+\norm{\psi''}_{L^\infty(I)}
  \label{eq:L-choice}
\end{equation}
and consider potentials of the form
\begin{equation}
  \Phi(x,y)
  =\psi(x)+\varepsilon h(x)y+\frac{L}{2}y^2+\frac{b(x)}6y^3,
  \label{eq:Phi}
\end{equation}
where $\varepsilon>0$ will be chosen small and the smooth function $b$ will
be fixed after that choice.  On $y=0$,
\begin{equation}
  D^2\Phi(x,0)
  =
  \begin{pmatrix}
    \psi''(x)&\varepsilon g(x)\\
    \varepsilon g(x)&L
  \end{pmatrix}.
  \label{eq:Hessian-axis}
\end{equation}
Set
\begin{equation}
  d(x):=L-\psi''(x),
  \qquad
  d_0:=\min_{x\in I}d(x)\ge2.
  \label{eq:d-d0}
\end{equation}
The two eigenvalues of \cref{eq:Hessian-axis} are
\begin{equation}
  \lambda_{1,2}(x,0)
  =\frac{\psi''(x)+L
  \mp\sqrt{d(x)^2+4\varepsilon^2g(x)^2}}{2}.
  \label{eq:eigenvalues-axis}
\end{equation}
Their gap is therefore
\begin{equation}
  \lambda_2(x,0)-\lambda_1(x,0)
  =\sqrt{d(x)^2+4\varepsilon^2g(x)^2}
  \ge d_0\ge2.
  \label{eq:eigenvalue-gap-axis}
\end{equation}

We next fix orientations of the two eigenvectors.  Define
\begin{equation}
  \Delta(x):=\sqrt{d(x)^2+4\varepsilon^2g(x)^2},
  \qquad
  \tau(x):=-\frac{2\varepsilon g(x)}{\Delta(x)+d(x)},
  \label{eq:Delta-tau}
\end{equation}
and
\begin{equation}
  c(x):=\frac{1}{\sqrt{1+\tau(x)^2}},
  \qquad
  \sigma(x):=\tau(x)c(x).
  \label{eq:c-sigma}
\end{equation}
Then
\begin{equation}
  r_1(x,0):=(c(x),\sigma(x)),
  \qquad
  r_2(x,0):=(-\sigma(x),c(x))
  \label{eq:eigenvectors-axis}
\end{equation}
are unit eigenvectors corresponding to the lower and upper eigenvalues in
\cref{eq:eigenvalues-axis}.  Indeed,
\begin{equation}
  \frac{\lambda_1(x,0)-\psi''(x)}{\varepsilon g(x)}
  =-\frac{2\varepsilon g(x)}{\Delta(x)+d(x)}=\tau(x)
\end{equation}
when $g(x)\neq0$, and the formula extends continuously across the zeros of
$g$.  The second vector is the positively oriented orthogonal complement of
the first.

Because $\Delta\ge d>0$,
\begin{equation}
  \abs{\tau(x)}
  =\frac{2\varepsilon\abs{g(x)}}{\Delta(x)+d(x)}
  \le\frac{\varepsilon\abs{g(x)}}{d(x)}
  \le\frac{\varepsilon\norm{g}_{L^\infty(I)}}{d_0}.
  \label{eq:tau-bound}
\end{equation}
Moreover, $\abs{\tau(x)}<1$, so
\begin{equation}
  c(x)>\frac1{\sqrt2},
  \qquad
  \operatorname{sgn}\sigma(x)=-\operatorname{sgn}g(x),
  \label{eq:c-sigma-sign}
\end{equation}
where the sign equality is understood to give $\sigma=0$ when $g=0$.
It follows from \cref{eq:tau-bound} that, uniformly on $I$,
\begin{equation}
  \sigma\longrightarrow0,
  \qquad
  c\longrightarrow1
  \qquad\text{as }\varepsilon\longrightarrow0.
  \label{eq:eigenvector-limit}
\end{equation}

Along $y=0$, the only possibly nonzero components of $D^3\Phi$ are
\begin{equation}
  \Phi_{xxx}=\psi''',
  \qquad
  \Phi_{xxy}=\varepsilon g',
  \qquad
  \Phi_{xyy}=0,
  \qquad
  \Phi_{yyy}=b.
  \label{eq:third-derivatives-axis}
\end{equation}
Consequently,
\begin{align}
  D^3\Phi[r_1,r_1,r_1]
  &=A_1+b\sigma^3,
  \label{eq:first-contraction}\\
  D^3\Phi[r_2,r_2,r_2]
  &=A_2+bc^3,
  \label{eq:second-contraction}
\end{align}
where all quantities in \cref{eq:first-contraction,eq:second-contraction} are
evaluated at $(x,0)$ and
\begin{align}
  A_1(x)&:=\psi'''(x)c(x)^3
  +3\varepsilon g'(x)c(x)^2\sigma(x),
  \label{eq:A1}\\
  A_2(x)&:=-\psi'''(x)\sigma(x)^3
  +3\varepsilon g'(x)\sigma(x)^2c(x).
  \label{eq:A2}
\end{align}
The uniform limits in \cref{eq:eigenvector-limit} imply
\begin{equation}
  A_1+\sigma^3\longrightarrow\psi'''
  \quad\text{uniformly on }I,
  \qquad
  A_2+c^3\longrightarrow1
  \quad\text{uniformly on }I.
  \label{eq:A-limits}
\end{equation}
Choose $\varepsilon>0$ sufficiently small that
\begin{equation}
  A_1(x)+\sigma(x)^3\ge\frac\mu2
  \quad\text{for every }x\in E,
  \label{eq:epsilon-choice-first}
\end{equation}
and
\begin{equation}
  A_2(x)+c(x)^3\ge\frac12
  \quad\text{for every }x\in I.
  \label{eq:epsilon-choice-second}
\end{equation}
Such a choice is possible by \cref{eq:mu,eq:A-limits}.

It remains to choose $b$.  Since $g\ge0$ on $E$, we have
$\sigma\le0$ there by \cref{eq:c-sigma-sign}.  Thus
\cref{eq:epsilon-choice-first} implies
\begin{equation}
  A_1\ge\frac\mu2-\sigma^3>0
  \qquad\text{on }E.
  \label{eq:A1-positive-E}
\end{equation}
Define the bad set
\begin{equation}
  \mathcal B:=\set{x\in I:A_1(x)\le0}.
  \label{eq:bad-set}
\end{equation}
It is a closed subset of the compact interval $I$, hence compact, and
\cref{eq:A1-positive-E} shows that
\begin{equation}
  \mathcal B\Subset\set{x\in\R:g(x)<0}.
  \label{eq:bad-set-containment}
\end{equation}
The notation $\Subset$ in \cref{eq:bad-set-containment} is justified as
follows: $\mathcal B$ is compact, the set $\{g<0\}$ is open, and
$\mathcal B\cap\{g\ge0\}=\varnothing$.

If $\mathcal B=\varnothing$, take $\chi=0$ and $M=0$.  If
$\mathcal B\neq\varnothing$, the smooth cutoff lemma gives a function
\begin{equation}
  \chi\in C_c^\infty(\set{x\in\R:g(x)<0}),
  \qquad
  0\le\chi\le1,
  \qquad
  \chi=1\ \text{on }\mathcal B.
  \label{eq:chi}
\end{equation}
On $\mathcal B$ one has $g<0$ and hence $\sigma>0$.  Compactness gives
$\min_{\mathcal B}\sigma>0$.  Choose
\begin{equation}
  M>
  \max_{x\in\mathcal B}
  \left(\frac{-A_1(x)-\sigma(x)^3}{\sigma(x)^3}\right)_+
  \label{eq:M-choice}
\end{equation}
and define
\begin{equation}
  b(x):=1+M\chi(x).
  \label{eq:b-definition}
\end{equation}
Here $z_+:=\max\{z,0\}$.  In the case $\mathcal B=\varnothing$, condition
\cref{eq:M-choice} is omitted and the choice $M=0$ is retained.

We now verify the first cubic contraction on all of $I$.  There are three
cases.
\begin{enumerate}
  \item If $g(x)\ge0$, then $x\in E$ and $\chi(x)=0$, because
  $\supp\chi\subset\{g<0\}$.  Hence $b(x)=1$ and
  \begin{equation}
    A_1(x)+b(x)\sigma(x)^3
    =A_1(x)+\sigma(x)^3
    \ge\frac\mu2>0.
  \end{equation}
  \item If $g(x)<0$ and $A_1(x)>0$, then $\sigma(x)>0$ and $b(x)\ge1$, so
  \begin{equation}
    A_1(x)+b(x)\sigma(x)^3>0.
  \end{equation}
  \item If $g(x)<0$ and $A_1(x)\le0$, then $x\in\mathcal B$,
  $\chi(x)=1$, and the choice \cref{eq:M-choice} gives
  \begin{equation}
    A_1(x)+(1+M)\sigma(x)^3>0.
  \end{equation}
\end{enumerate}
These cases exhaust $I$, and therefore
\begin{equation}
  D^3\Phi(x,0)[r_1(x,0),r_1(x,0),r_1(x,0)]>0,
  \qquad x\in I.
  \label{eq:first-positive-axis}
\end{equation}
For the second contraction, \cref{eq:epsilon-choice-second} and $M\chi c^3
\ge0$ give directly
\begin{align}
  D^3\Phi(x,0)[r_2(x,0),r_2(x,0),r_2(x,0)]
  &=(A_2+c^3)+M\chi c^3\notag\\
  &\ge\frac12>0,
  \qquad x\in I.
  \label{eq:second-positive-axis}
\end{align}

\subsection{Passage to a convex open strip}

The preceding argument gives strict inequalities only on the central segment.
We now justify in detail that they persist on a convex open neighborhood.  A
direct differentiation of \cref{eq:Phi} gives the full Hessian
\begin{equation}
  D^2\Phi(x,y)
  =
  \begin{pmatrix}
    \psi''(x)+\varepsilon g'(x)y+\frac16b''(x)y^3
    &\varepsilon g(x)+\frac12b'(x)y^2\\[1mm]
    \varepsilon g(x)+\frac12b'(x)y^2
    &L+b(x)y
  \end{pmatrix}.
  \label{eq:full-Hessian}
\end{equation}
In particular,
\begin{equation}
  \Phi_{yy}(x,y)-\Phi_{xx}(x,y)
  =d(x)+\bigl(b(x)-\varepsilon g'(x)\bigr)y
  -\frac16b''(x)y^3.
  \label{eq:diagonal-difference}
\end{equation}
The functions in \cref{eq:diagonal-difference} are smooth and bounded on a
neighborhood of $I$.  In view of \cref{eq:d-d0}, there is $\delta_1>0$ such
that
\begin{equation}
  \Phi_{yy}(x,y)-\Phi_{xx}(x,y)\ge\frac{d_0}{2}>0
\end{equation}
for $(x,y)\in I\times[-\delta_1,\delta_1]$.

For such $(x,y)$, set
\begin{equation}
  \mathfrak d(x,y):=\Phi_{yy}(x,y)-\Phi_{xx}(x,y),
  \qquad
  \mathfrak z(x,y):=\Phi_{xy}(x,y),
\end{equation}
\begin{equation}
  \widehat\Delta(x,y)
  :=\sqrt{\mathfrak d(x,y)^2+4\mathfrak z(x,y)^2},
  \qquad
  \widehat\tau(x,y)
  :=-\frac{2\mathfrak z(x,y)}{
  \widehat\Delta(x,y)+\mathfrak d(x,y)},
  \label{eq:off-axis-tau}
\end{equation}
and
\begin{equation}
  \widehat c(x,y):=(1+\widehat\tau(x,y)^2)^{-1/2},
  \qquad
  \widehat\sigma(x,y):=\widehat\tau(x,y)\widehat c(x,y).
\end{equation}
Then
\begin{equation}
  \widehat r_1(x,y):=(\widehat c,\widehat\sigma),
  \qquad
  \widehat r_2(x,y):=(-\widehat\sigma,\widehat c)
  \label{eq:off-axis-eigenvectors}
\end{equation}
are smooth, consistently oriented unit eigenvectors of $D^2\Phi(x,y)$ for its
lower and upper eigenvalues.  Their eigenvalue gap is
\begin{equation}
  \widehat\lambda_2(x,y)-\widehat\lambda_1(x,y)
  =\widehat\Delta(x,y)
  \ge\mathfrak d(x,y)
  \ge\frac{d_0}{2}.
  \label{eq:off-axis-gap}
\end{equation}
Thus the Hessian has two distinct real eigenvalues on the closed strip of
height $2\delta_1$.  At $y=0$, the definitions in
\cref{eq:off-axis-tau,eq:off-axis-eigenvectors} reduce to
\cref{eq:Delta-tau,eq:c-sigma,eq:eigenvectors-axis}; in particular,
$\widehat r_k(x,0)=r_k(x,0)$.

For $k=1,2$, define
\begin{equation}
  \Psi_k(x,y)
  :=D^3\Phi(x,y)
  [\widehat r_k(x,y),\widehat r_k(x,y),\widehat r_k(x,y)].
  \label{eq:Psi-k}
\end{equation}
These functions are continuous.  By
\cref{eq:first-positive-axis,eq:second-positive-axis}, they are strictly
positive on the compact segment $I\times\{0\}$.  Hence
\begin{equation}
  m_k:=\min_{x\in I}\Psi_k(x,0)>0,
  \qquad k=1,2.
\end{equation}
Uniform continuity on a compact rectangle permits us to choose
$0<\delta\le\delta_1$ so small that
\begin{equation}
  \Psi_k(x,y)\ge\frac{m_k}{2}>0,
  \qquad
  (x,y)\in I\times[-\delta,\delta],
  \qquad k=1,2.
  \label{eq:strip-positive}
\end{equation}
We now fix this $\delta$ and take $\V$ as in \cref{eq:state-domain}.  Since
$DF=D^2\Phi$, \cref{eq:off-axis-gap} proves strict hyperbolicity of the flux
$F=\nabla\Phi$ on $\V$.  Combining \cref{lem:eigenvalue-derivative} with
\cref{eq:strip-positive} gives
\begin{equation}
  D\lambda_k(U)[\widehat r_k(U)]
  =D^3\Phi(U)[\widehat r_k(U),\widehat r_k(U),\widehat r_k(U)]
  >0,
  \qquad U\in\V,\qquad k=1,2.
  \label{eq:genuine-nonlinearity-final}
\end{equation}
Thus both characteristic fields are genuinely nonlinear throughout the same
convex open state domain.

\section{Entropy compatibility and recovery of the four wells}
\label{sec:recovery}

With the choices above, define on $\V$
\begin{equation}
  F:=\nabla\Phi,
  \qquad
  \eta(x,y):=\frac{x^2+y^2}{2},
  \qquad
  q(x,y):=(x,y)\cdot\nabla\Phi(x,y)-\Phi(x,y).
  \label{eq:final-F-eta-q}
\end{equation}
Equivalently, the two components of the flux and the entropy flux are
\begin{align}
  F_1(x,y)
  &=\psi'(x)+\varepsilon g(x)y+\frac16b'(x)y^3,
  \label{eq:explicit-F1}\\
  F_2(x,y)
  &=\varepsilon h(x)+Ly+\frac12b(x)y^2,
  \label{eq:explicit-F2}\\
  q(x,y)
  &=x\psi'(x)-\psi(x)+\varepsilon xg(x)y+\frac{L}{2}y^2
    +\frac16xb'(x)y^3+\frac13b(x)y^3.
  \label{eq:explicit-q}
\end{align}
Indeed, the first two identities follow by differentiating
\cref{eq:Phi}, and the third follows by substituting them into
$q=(x,y)\cdot F-\Phi$.
All three functions are smooth on a neighborhood of the closed rectangle
$I\times[-\delta,\delta]$.  By \cref{lem:potential-entropy},
\begin{equation}
  D^2\eta=I_2>0,
  \qquad
  Dq=D\eta\,DF
  \quad\text{on }\V.
  \label{eq:final-entropy-check}
\end{equation}

At $U_i=(\xi_i,0)$, \cref{eq:h-zeros,eq:Phi} give
\begin{equation}
  \Phi(U_i)=\psi(\xi_i),
  \qquad
  \Phi_x(U_i)=\psi'(\xi_i),
  \qquad
  \Phi_y(U_i)=\varepsilon h(\xi_i)=0.
  \label{eq:jets-at-Ui}
\end{equation}
Therefore
\begin{equation}
  F(U_i)=\bigl(\psi'(\xi_i),0\bigr),
  \qquad
  q(U_i)=\xi_i\psi'(\xi_i)-\psi(\xi_i).
  \label{eq:F-q-at-Ui}
\end{equation}
Substituting the exact data from \cref{eq:Hermite-verification} yields
\begin{equation}
\renewcommand{\arraystretch}{1.25}
\begin{array}{c|rrrr}
  i&1&2&3&4\\
  \hline
  \xi_i&-7&-1&1&7\\
  -F_1(U_i)&-\frac{52}{5}&-\frac{34}{5}&6&0\\[1mm]
  -F_2(U_i)&0&0&0&0\\[1mm]
  \eta(U_i)&\frac{49}{2}&\frac12&\frac12&\frac{49}{2}\\[1mm]
  -q(U_i)&\frac{48}{5}&\frac{216}{5}&-24&0.
\end{array}
\label{eq:well-data-table}
\end{equation}
Combining \cref{eq:G-map,eq:well-data-table} gives, entry by entry,
\begin{equation}
  G_{F,\eta,q}(U_i)=X_i,
  \qquad i=1,2,3,4.
  \label{eq:wells-recovered}
\end{equation}
All four states lie in the interior of $\V$ because
$\xi_i\in(-8,8)$ and $0\in(-\delta,\delta)$.  It follows that
\begin{equation}
  (X_1,X_2,X_3,X_4)
  \in\Kprod_{F,\eta,q}^{(4)}.
\end{equation}
The $T_4$ identities were verified in \cref{sec:algebraic-T4}, and
\cref{eq:rank-two-conclusion} shows that the tuple is nondegenerate.  Strict
hyperbolicity and genuine nonlinearity were proved in
\cref{eq:off-axis-gap,eq:genuine-nonlinearity-final}, while
\cref{eq:final-entropy-check} supplies the strictly convex entropy and its
compatible entropy flux.  This completes the proof of \cref{thm:main}.
\qed

\appendix

\section{Exact-arithmetic checks}
\label{app:exact-checks}

This appendix records compact checks of the numerical identities used in the
construction.  First, the rank-one factorizations and the closure condition
can be checked entirely at the level of the two column vectors:
\begin{align}
  &\begin{pmatrix}-26\\0\\24\end{pmatrix}
  +\begin{pmatrix}2\\0\\-48\end{pmatrix}
  +\begin{pmatrix}6\\0\\-24\end{pmatrix}
  +\begin{pmatrix}18\\0\\48\end{pmatrix}
  =0,\\
  &\begin{pmatrix}-26\\0\\24\end{pmatrix}
  +\begin{pmatrix}-4\\0\\96\end{pmatrix}
  +\begin{pmatrix}30\\0\\-120\end{pmatrix}
  +\begin{pmatrix}0\\0\\0\end{pmatrix}
  =0,
\end{align}
after multiplying the matrices by $5$.  Second, the correction polynomial in
\cref{eq:psi-definition} may be expanded as
\begin{equation}
  xR(x)^2=x^9-100x^7+2598x^5-4900x^3+2401x.
  \label{eq:correction-expanded}
\end{equation}
The vanishing of this polynomial and its first derivative at
$x=\pm1,\pm7$ follows either from \cref{eq:R-polynomial} or directly from
\cref{eq:correction-expanded}.  Substitution into \cref{eq:psi-H} gives the
eight rational values in \cref{eq:Hermite-verification}; differentiating
three times and substituting $x=-5,0,5$ gives exactly the three fractions in
\cref{eq:positive-third-derivatives}.

Finally, the lower-right entries in the four wells are recovered from
\begin{equation}
  \psi(\xi_i)-\xi_i\psi'(\xi_i)
  =\left(\frac{48}{5},\frac{216}{5},-24,0\right)_i,
\end{equation}
whereas their upper-right entries are
\begin{equation}
  -\psi'(\xi_i)
  =\left(-\frac{52}{5},-\frac{34}{5},6,0\right)_i.
\end{equation}
These are precisely the second-column entries in
\cref{eq:X12,eq:X34}.

\section{The smooth cutoff fact used above}
\label{app:cutoff}

We used the following elementary one-dimensional form of the smooth cutoff
lemma.  If $K\subset\Omega\subset\R$, where $K$ is compact and $\Omega$ is
open, then there exists $\chi\in C_c^\infty(\Omega)$ such that
$0\le\chi\le1$ and $\chi=1$ on $K$.  To see this, choose open sets
$O_1,O_2$ satisfying
\begin{equation}
  K\subset O_1,
  \qquad
  \overline{O_1}\subset O_2,
  \qquad
  \overline{O_2}\subset\Omega,
\end{equation}
which is possible by compactness.  Convolve the indicator function of $O_2$
with a nonnegative standard mollifier whose support is smaller than both
$\operatorname{dist}(\overline{O_1},\R\setminus O_2)$ and
$\operatorname{dist}(\overline{O_2},\R\setminus\Omega)$.  The resulting
smooth function is supported in $\Omega$, takes values in $[0,1]$, and equals
one on $O_1$, hence on $K$.

\end{document}
